\textbf{(i)} The prime correspondence identifies $\operatorname{Spec}(S^{-1}A)$ with the primes of $A$ disjoint from $S$. Under this bijection, $D(a/s)$ maps to $D(a)\cap S^{-1}X$, so the bijection and its inverse are continuous. For $S=\{1,a,a^2,\ldots\}$, disjointness is equivalent to not containing $a$.

\textbf{(ii)} A prime $\mathfrak q\subseteq B$ is disjoint from $f(S)$ exactly when $f^{-1}(\mathfrak q)$ is disjoint from $S$. Contraction also commutes with localization, so the maps agree on these subspaces.

\textbf{(iii)} The quotient correspondences identify the two spectra with $V(\mathfrak a)$ and $V(\mathfrak b)$. A prime $\mathfrak q/\mathfrak b$ contracts to $f^{-1}(\mathfrak q)/\mathfrak a$, proving the assertion.

\textbf{(iv)} Take $S=A\setminus\mathfrak p$. By (ii), the primes of $B_{\mathfrak p}$ correspond to primes of $B$ whose contractions lie in $\mathfrak p$. By (iii), quotienting by $\mathfrak pB_{\mathfrak p}$ retains exactly those whose contractions contain $\mathfrak p$, hence equal $\mathfrak p$. Both identifications are homeomorphisms onto the relevant subspaces. Finally, $B_{\mathfrak p}/\mathfrak pB_{\mathfrak p}\cong k(\mathfrak p)\otimes_A B$.
